\documentclass[conference]{IEEEtran}
\IEEEoverridecommandlockouts
% The preceding line is only needed to identify funding in the first footnote. If that is unneeded, please comment it out.
\usepackage{cite}
\usepackage{amsmath,amssymb,amsfonts}
\usepackage{algorithmic}
\usepackage{graphicx}
\usepackage{textcomp}
\usepackage{xcolor}
\def\BibTeX{{\rm B\kern-.05em{\sc i\kern-.025em b}\kern-.08em
    T\kern-.1667em\lower.7ex\hbox{E}\kern-.125emX}}
\begin{document}

\title{Latent Boundary Detection for Rollout Re-Anchoring in Frozen JEPA World Models}

\author{\IEEEauthorblockN{Jiaqi Wang}
\IEEEauthorblockA{\textit{School of Mechanical and Electrical Engineering} \\
\textit{Guangzhou University}\\
Guangzhou, China}
\and
\IEEEauthorblockN{Tao Zou\textsuperscript{*}}
\IEEEauthorblockA{\textit{School of Mechanical and Electrical Engineering} \\
\textit{Guangzhou University}\\
Guangzhou, China \\
\textsuperscript{*}Corresponding author}
}

\maketitle

\begin{abstract}
% TODO: 摘要最后写。一句话主张：无监督变点切分在冻结JEPA潜空间自动发现物理边界，
% 蒸馏成轻量在线检测头，边界触发停止重锚提升长距离规划成功率。
\end{abstract}

\begin{IEEEkeywords}
artificial intelligence, world models,  representation learning
\end{IEEEkeywords}

\section{Introduction}
An agent that acts in the world benefits from predicting the consequences of its actions, and
world models are learned for this purpose~\cite{lecun2022path}. Joint-Embedding Predictive
Architectures (JEPA) are a family of world models that predict latent embeddings instead of raw
pixels, which avoids pixel reconstruction and keeps the model lightweight. Such models have
been applied to image and video understanding (I-JEPA~\cite{assran2023ijepa}, V-JEPA~\cite{vjepa2024})
and to robot control, where LeWM pairs the predictive model with a sampling-based planner
that rolls out candidate action sequences in imagination and executes the best one, as in
model predictive control~\cite{lewm2026}. Yet rollouts of a frozen world model drift from the
real dynamics, and the error grows with every imagined step, so planning success falls with
task distance on PushT~\cite{mindthegap2026}.

To solve this problem, hierarchical planners such as H-JEPA plan at multiple timescales, but
an extra model must be trained end-to-end on top of the world model~\cite{hjepa2026}. Subgoal
planners shorten the effective horizon, but the subgoals they generate can be spurious and can
lower success~\cite{mindthegap2026}. A question remains: can boundary information replace
subgoals in keeping the rollout error small? If it can, knowing when to stop a rollout and
when to re-anchor it is enough.

To answer this question, a latent boundary detector is built. Three generic signals of the
latent stream are measured. Increment energy reflects the magnitude of the latent change.
Direction change reflects turns of the latent trajectory. Prediction novelty reflects the
mismatch between the predicted next state and the observed one. These signals are segmented
offline with an unsupervised changepoint method~\cite{killick2012pelt}, and the segmentation
is distilled into a lightweight causal head. During planning, the head
reads the imagined rollout stream and the real observation stream. A boundary event stops the
current plan and forces an immediate re-anchoring. The contributions are:
\begin{itemize}
\item A paired-task study verifying that the timing of stopping and re-anchoring controls
planning success in frozen world models.
\item A latent boundary detector (${\sim}3$k parameters) that decides when to stop a
rollout and re-anchor it, without task rules and without retraining the world model.
\item An evaluation on PushT, plus a rhythm-level diagnosis on OGB Cube that delimits
when boundary-triggered re-anchoring can help and when it cannot.
\end{itemize}


\section{Related Work}
% 每篇论文按统一逻辑写：是什么 → 解决了什么 → 发现了什么/剩下什么问题。
% 收紧版：intro 已写过的背景不再复述，只写与我们的差异。
\textbf{World models and JEPA.} LeCun's path framework frames intelligence as predictive world
models~\cite{lecun2022path}. I-JEPA showed that predicting latent embeddings rather than pixels
yields strong representations~\cite{assran2023ijepa}, and V-JEPA extended this recipe to video,
demonstrating that the world model itself need never decode pixels~\cite{vjepa2024}. LeWM
brought JEPA to robot control: a frozen latent predictor paired with sampling-based planning
rolls out actions in imagination and executes them in the real world~\cite{lewm2026}. Rollouts
of such a frozen model, however, drift from reality, and the drift compounds with every
imagined step~\cite{mindthegap2026}. Existing remedies train extra modules on top of the world
model, but retraining is costly: H-JEPA plans at multiple timescales and requires an extra
hierarchical model trained end-to-end~\cite{hjepa2026}. Subgoal-based planners shorten the
effective horizon by planning toward intermediate goals, but the subgoals they generate can be
spurious and lower success~\cite{mindthegap2026}. Our method requires neither: it reuses the
frozen world model as-is and reads the boundary signal from its own latent stream.

\textbf{Change-point detection.} Classical detectors segment time series at changes in
likelihood: CUSUM accumulates evidence for a single shift~\cite{page1954cusum}, PELT finds the
optimal multiple-change segmentation in linear cost~\cite{killick2012pelt}, and BOCPD gives a
Bayesian online treatment~\cite{adams2007bocpd}. These methods are our baselines: they capture
the offline signal, but they carry heavy distributional assumptions and deliver
no lightweight online head for the planner. We distill their boundary structure into a
$\sim$3k-parameter causal CNN that runs on the latent stream.

\textbf{Action chunking and adaptive execution horizons.} Chunking regresses multiple actions
at once to tame error compounding~\cite{act2023}. Fixed chunk lengths, however, are
task-specific: they overcommit in quiet stretches and undercommit at decisive events. AutoHorizon
estimates execution horizons dynamically from self-attention weights of a VLA policy at test
time~\cite{autohorizon2026}; PACE identifies low-speed transitions in the predicted action
chunk and re-anchors at those phase boundaries, without retraining~\cite{pace2026}. Both are
test-time adaptive-horizon methods, and they are our closest relatives: our method resembles a
dynamic horizon $h$. The difference is where the dynamics live. We adapt horizons \emph{inside the
planner}: the detector sits on the frozen world model's latent stream, sets the planning
horizon $h$ from the imagined fire step, and truncates the execution walk on the real stream to
$\min(h,k)$. Our boundaries come from generic latent signals (increment energy, direction
change, prediction novelty), not from policy internals or action-space heuristics, and the
detector never enters the planner's search loop. This lineage goes back to event-triggered
control, where a controller recomputes the moment an event fires rather than on a fixed
clock~\cite{heemels2012}; our imagined and real posts are two such triggers on the same latent
stream, one in simulation and one in reality.

\section{Method}

\subsection{Problem Setting}
A frozen LeWM encodes pixels into latent states $\mathbf{z}_t\in\mathbb{R}^d$ ($d=192$; one model step covers $m=5$ frames). Actions are planned in blocks of $B=5$ steps. For a plan horizon of $h$ blocks ($T=hB$ steps), CEM searches the action sequence that minimizes goal error evaluated only at the rollout end:
\begin{equation}
\min_{\mathbf{a}_{1:T}}\ \lVert \hat{\mathbf{z}}_{t+T}-\mathbf{g}\rVert_2^2 ,
\end{equation}
where $\hat{\mathbf{z}}$ is the imagined state and $\mathbf{g}$ the goal embedding. The executed plan is re-anchored after $r$ blocks. A frozen model drifts as the rollout lengthens, so the schedule of $h$ and $r$ is exactly what this paper controls.

\subsection{Latent Boundary Signals}
Three generic signals are read from the latent stream (model step $t$):
\begin{align}
f_1[t] &= \lVert \mathbf{z}_t-\mathbf{z}_{t-1}\rVert, \\
f_2[t] &= \frac{\langle\Delta\mathbf{z}_t,\Delta\mathbf{z}_{t-1}\rangle}{\lVert\Delta\mathbf{z}_t\rVert\,\lVert\Delta\mathbf{z}_{t-1}\rVert}, \\
f_3[t] &= \lVert \mathbf{z}_{t+1}-F_z(\mathbf{z}_{t-2:t})\rVert,
\end{align}
with $\Delta\mathbf{z}_t=\mathbf{z}_t-\mathbf{z}_{t-1}$. $F_z{:}\ \mathbb{R}^{3d}\to\mathbb{R}^d$ is a linear forward model fitted by least squares on training streams and is action-blind by design: it receives no actions, so action-driven kinks stay surprising to it even inside imagined rollouts. The three signals look at the same stream through different lenses: $f_1$ sees contact energy, $f_2$ sees reversals of the latent trajectory, $f_3$ sees surprise against a fixed predictor.

\subsection{Detector Head}
For each signal, causal statistics over windows $W\in\{2,4,8\}$ give 18 channels, $\phi_t=\{\mu_W(f_i),\sigma_W(f_i)\}_{i,W}$, z-scored with training statistics. A causal CNN (${\sim}3$k parameters) reads only past rows and outputs the boundary probability $p_t=\sigma(\cdot)$. Training uses binary labels of physical contact events derived from kinematic state thresholds (24.0\% positive, weight 3.16) with hold-AUC early stopping; the offline PELT segmentation~\cite{killick2012pelt} is kept as an unsupervised baseline. At inference, a fire is declared when $p_t>\tau$ and the refractory $\rho=5$ has elapsed since the last fire ($\tau=0.5$).

\subsection{Dual-Post Re-Anchoring}
The same head occupies two posts, and both are active during a rollout.
\textit{Planning (imagined) post.} Before each CEM query, one probe rollout extends the latest 3 real latents by $H=10$ zero-action steps in imagination. The first fire index $s$ (in model steps) on this imagined sequence sets the horizon directly, $h=\max(s,5)$; if nothing fires, $h=10$, the full probe length. The head never enters the CEM loop, so candidates cannot game it, and warm start is disabled.
\textit{Execution (real) post.} While a plan executes, the head also scores each incoming real latent, using a buffer of up to 12 frames. A fire discards the unexecuted buffer and forces an immediate re-anchor, so the walk is truncated to $\min(h,k)$, where $k$ is the distance to that real-stream fire. Both lengths are therefore decided by the detector: one in imagination, one in reality.

\section{Experiments}
% TODO: E1 复现基线 / Table 1 消融+对齐 / Table 2 三臂成功率 / Fig 2 漂移清零
% 口径（2026-11-05 用户裁决）：confirmatory 主线 = PushT（Task A 决定性三角 n=150 + Task B b2 argmax 网格）；
% cube 节律谱系关卡 FAIL 作为"条件性判据"写入 Limitation/条件分析段，不做评测管线。禁"两任务正结果"表述
\section{Conclusion}
% TODO: 一段总结 + limitation（单环境、蒸馏误差）

% ---- 图表占位（写正文时插入并引用）----
% Fig. 1: 管线图 潜变量流 -> 三路特征 -> PELT切分 -> student头 -> 边界停止重锚
% Fig. 2: 一条轨迹重锚定前后 rollout 漂移曲线（漂移清零）
% Table I: 检测消融 + 物理事件对齐 F1（特征递增 / 切分方法 / 头部）
% Table II: 三臂成功率主表（平坦 vs 边界停止 vs 等频率停止, d=25/50/75）

% 手动管理：条目顺序 = 正文首次引用顺序（IEEE 规范）。
% 后续每新增一条引用，把 \bibitem 追加到列表末尾，不要插到中间。
\begin{thebibliography}{00}

\bibitem{lecun2022path} Y. LeCun, ``A path towards autonomous machine intelligence,'' OpenReview, version 0.9.2, 2022.

\bibitem{assran2023ijepa} M. Assran, Q. Duval, I. Misra, P. Bojanowski, P. Vincent, M. Rabbat, Y. LeCun, and N. Ballas, ``Self-supervised learning from images with a joint-embedding predictive architecture,'' in Proc. IEEE/CVF Conf. Comput. Vis. Pattern Recognit. (CVPR), 2023, pp. 15619--15629.

\bibitem{vjepa2024} A. Bardes, Q. Garrido, J. Ponce, X. Chen, M. Rabbat, Y. LeCun, M. Assran, and N. Ballas, ``Revisiting feature prediction for learning visual representations from video,'' in Proc. Int. Conf. Mach. Learn. (ICML), 2024.

\bibitem{lewm2026} L. Maes, Q. Le Lidec, D. Scieur, Y. LeCun, and R. Balestriero, ``LeWorldModel: Stable end-to-end joint-embedding predictive architecture from pixels,'' arXiv preprint arXiv:2603.19312, 2026.

\bibitem{hjepa2026} W. Zhang, B. Terver, M. Rabbat, Y. LeCun, and R. Balestriero, ``H-JEPA: End-to-end learning of hierarchical world models for visual planning,'' arXiv preprint arXiv:2610.06805, 2026.

\bibitem{mindthegap2026} N. Caselli, F. Massafra, S. Punzo, S. Lo Sardo, I. Pantelidis, and S. K. Bhethanabhotla, ``Mind the gap: Promises and pitfalls of hierarchical planning in LeWorldModel,'' arXiv preprint arXiv:2607.12547, 2026.

\bibitem{killick2012pelt} R. Killick, P. Fearnhead, and I. A. Eckley, ``Optimal detection of changepoints with a linear computational cost,'' J. Amer. Statist. Assoc., vol. 107, no. 500, pp. 1590--1598, 2012.

\bibitem{page1954cusum} E. S. Page, ``Continuous inspection schemes,'' Biometrika, vol. 41, no. 1/2, pp. 100--115, 1954.

\bibitem{adams2007bocpd} R. P. Adams and D. J. C. MacKay, ``Bayesian online changepoint detection,'' arXiv preprint arXiv:0710.3742, 2007.

\bibitem{act2023} T. Z. Zhao, V. Kumar, S. Levine, and C. Finn, ``Learning fine-grained bimanual manipulation with low-cost hardware,'' in Proc. Robot.: Sci. Syst. (RSS), 2023.

\bibitem{autohorizon2026} H. Wang, G. Zhang, Y. Yan, R. R. Kompella, and G. Liu, ``VLA knows its limits: Adaptive execution horizons for robot policies,'' in Proc. Eur. Conf. Comput. Vis. (ECCV), 2026.

\bibitem{pace2026} J. Nie, J. Li, C. Liu, J. Lao, J. Zhang, T. Zhang, L. Lin, and S. Huang, ``PACE: Phase-aware chunk execution for robot policies with action chunking,'' arXiv preprint arXiv:2606.00537, 2026.

\bibitem{heemels2012} W. P. M. H. Heemels, K. H. Johansson, and P. Tabuada, ``An introduction to event-triggered and self-triggered control,'' in Proc. IEEE Conf. Decis. Control (CDC), 2012, pp. 3270--3285.

% TODO: 核实后加入 Basseville & Nikiforov 1993 (Detection of Abrupt Changes)

\end{thebibliography}

\end{document}
